% !TEX root = ../main.tex
\section{Cơ sở lý thuyết}

Thuật toán Di truyền (Genetic Algorithm - GA) là một phương pháp tối ưu hóa ngẫu nhiên được lấy cảm hứng từ quy luật tiến hóa và di truyền trong tự nhiên. Quy trình cốt lõi của GA dựa trên việc duy trì và biến đổi một quần thể các cá thể, trong đó mỗi cá thể đại diện cho một nghiệm tiềm năng của bài toán. Quá trình tiến hóa bắt đầu bằng việc khởi tạo quần thể, sau đó trải qua các vòng lặp đánh giá độ thích nghi (fitness), chọn lọc các cá thể ưu tú, và áp dụng các toán tử lai ghép (crossover) cũng như đột biến (mutation) để sinh ra thế hệ kế tiếp với chất lượng ngày càng cải thiện.

Trong bài toán xếp lịch, việc mã hóa các cá thể đóng vai trò quyết định đến hiệu quả của quá trình tìm kiếm. Chúng tôi áp dụng phương pháp mã hóa trực tiếp (direct encoding), trong đó mỗi gene lưu trữ đầy đủ thông tin về một lớp học phần bao gồm ngày học, khung giờ và định danh phòng học. Cấu trúc này cho phép các toán tử di truyền tác động trực tiếp lên các thuộc tính vật lý của lịch học, giúp duy trì tính nhất quán và dễ dàng kiểm soát các ràng buộc giữa các buổi học trong cùng một môn học.

Hàm thích nghi (fitness function) được thiết kế nhằm phản ánh mức độ hoàn thiện của lịch học thông qua việc tổng hợp các hình phạt từ các vi phạm ràng buộc. Các ràng buộc trong hệ thống được phân cấp rõ rệt thành ràng buộc cứng và ràng buộc mềm. Ràng buộc cứng bao gồm các điều kiện bắt buộc phải thỏa mãn như tính khả dụng của phòng học, sự phù hợp về dung tích và việc tránh xung đột thời gian giữa các giảng viên hoặc phòng học. Ngược lại, các ràng buộc mềm hướng tới việc nâng cao chất lượng lịch học bằng cách tối ưu hóa sự phân bổ đều đặn các lớp học xuyên suốt tuần và giảm thiểu các khoảng thời gian trống giữa các tiết học của cùng một lớp.

Công thức tính độ thích nghi thường được biểu diễn dưới dạng nghịch đảo của tổng phạt trọng số, trong đó các vi phạm ràng buộc cứng được gán các hệ số phạt rất lớn để đảm bảo thuật toán ưu tiên tìm kiếm các lịch học hợp lệ trước khi tối ưu hóa các yếu tố chất lượng. Thông qua cơ chế chọn lọc, các cá thể có độ thích nghi cao sẽ có xác suất lớn hơn được giữ lại và truyền lại các đặc tính tốt cho thế hệ sau. Chiến lược bảo tồn cá thể ưu tú (elitism) cũng được tích hợp để đảm bảo rằng nghiệm tốt nhất tìm được qua các thế hệ luôn được duy trì, giúp thuật toán hội tụ một cách ổn định về phía nghiệm tối ưu toàn cục.
